\documentclass[11pt]{article}
\usepackage{amsmath,amssymb,amsthm,mathtools}
\usepackage[margin=1in]{geometry}
\usepackage{microtype}

\newtheorem{theorem}{Theorem}[section]
\newtheorem{lemma}[theorem]{Lemma}
\newtheorem{proposition}[theorem]{Proposition}
\newtheorem{corollary}[theorem]{Corollary}
\newtheorem{fact}[theorem]{Fact}
\theoremstyle{definition}
\newtheorem{definition}[theorem]{Definition}
\newtheorem{remark}[theorem]{Remark}
\newtheorem{conjecture}[theorem]{Conjecture}

\newcommand{\chiG}{\chi(G)}
\newcommand{\degen}{\operatorname{degen}}
\newcommand{\mad}{\operatorname{mad}}
\newcommand{\had}{\eta}

\title{Hadwiger's Conjecture:\\
Strongest Completely Elementary Affirmative Results}
\author{Self-contained synthesis from first principles}
\date{}

\begin{document}
\maketitle

\begin{abstract}
We prove, entirely from first principles and without black-box citations,
every affirmative result toward Hadwiger's conjecture that admits a fully
elementary proof in the classical repertoire:
\begin{enumerate}
\item Hadwiger's conjecture for all \(t\le 4\);
\item the Mader--Wagner bound \(\chi(G)\le 2^{t-2}\) for every \(K_t\)-minor-free
graph \(G\), together with the matching average-degree statement that every
graph of average degree at least \(2^{t-2}\) contains a (in fact, dominating)
\(K_t\)-model;
\item the standard reductions (components, blocks/cutvertices, criticality)
and the structural theory of minimal counterexamples / \(t\)-critical graphs
that can be established elementarily;
\item for \(t=5\), everything that does \emph{not} require the Four Colour
Theorem, together with a precise account of what does.
\end{enumerate}
No claim is made for general \(t\ge 6\) beyond the exponential bound. In
particular we do \emph{not} invoke Robertson--Seymour--Thomas, Kostochka--Thomason
density, or any computer-assisted argument.
\end{abstract}

\tableofcontents

%------------------------------------------------------------------------------
\section{Definitions and notation}
%------------------------------------------------------------------------------

Throughout, graphs are finite, undirected, and simple (no loops, no parallel
edges), unless multigraphs are explicitly allowed for intermediate contraction
steps. Write \(n(G)=|V(G)|\), \(e(G)=|E(G)|\), \(\delta(G)\) for minimum degree,
\(\Delta(G)\) for maximum degree, \(\kappa(G)\) for vertex-connectivity, and
\(d(G)=2e(G)/n(G)\) for average degree when \(n(G)>0\).

\begin{definition}[Models and minors]
A \emph{\(K_t\)-model} (or \emph{\(K_t\)-minor model}) in \(G\) is a family of
pairwise disjoint nonempty sets \(B_1,\dots,B_t\subseteq V(G)\) (``branch
sets'') such that
\begin{enumerate}
\item \(G[B_i]\) is connected for each \(i\);
\item for all \(i\neq j\) there is at least one edge of \(G\) with one end in
\(B_i\) and the other in \(B_j\).
\end{enumerate}
We write \(G\succcurlyeq K_t\) if such a model exists. Equivalently, \(K_t\) can
be obtained from a subgraph of \(G\) by contracting edges (and discarding
isolated vertices produced along the way). The \emph{Hadwiger number}
\(\had(G)\) is the largest \(t\) such that \(G\succcurlyeq K_t\).
\end{definition}

\begin{definition}[Dominating models]
A \emph{dominating \(K_t\)-model} in \(G\) is a sequence \((T_1,\dots,T_t)\) of
pairwise disjoint nonempty connected subgraphs of \(G\) such that for all
\(1\le i<j\le t\), every vertex of \(T_j\) has a neighbour in \(T_i\).
\end{definition}

Every dominating \(K_t\)-model is a \(K_t\)-model (take the vertex sets of the
\(T_i\)). The converse fails for \(t\ge 4\), but the two notions coincide for
\(t\le 3\) (Observation~\ref{obs:dom-small}).

\begin{definition}[Colouring and criticality]
A proper \(k\)-colouring is a map \(c:V(G)\to\{1,\dots,k\}\) with \(c(u)\neq c(v)\)
whenever \(uv\in E(G)\). The chromatic number \(\chi(G)\) is the least such \(k\).
The graph \(G\) is \emph{\(t\)-critical} if \(\chi(G)=t\) and \(\chi(H)<t\) for
every proper subgraph \(H\subsetneq G\).
\end{definition}

\begin{definition}[Degeneracy]
\(G\) is \emph{\(d\)-degenerate} if every nonempty subgraph has a vertex of
degree at most \(d\). Write \(\degen(G)\) for the least such \(d\), and
\(\mad(G)=\max\{d(H):H\subseteq G,\,H\neq\emptyset\}\).
\end{definition}

\begin{conjecture}[Hadwiger, 1943]
For every integer \(t\ge 1\) and every graph \(G\),
\[
G\not\succcurlyeq K_t \quad\Longrightarrow\quad \chi(G)\le t-1.
\]
Equivalently, \(\chi(G)\le\had(G)\) for every \(G\). Write \(\mathrm{HC}_t\) for
the case of \(K_t\)-minors (i.e.\ \(\chi\le t-1\) for \(K_t\)-minor-free graphs).
\end{conjecture}

%------------------------------------------------------------------------------
\section{Elementary facts on colouring, minors, and degeneracy}
%------------------------------------------------------------------------------

\begin{fact}[Monotonicity of minors]\label{fact:mono}
If \(H\) is a subgraph or minor of \(G\) and \(H\succcurlyeq K_t\), then
\(G\succcurlyeq K_t\). Consequently every subgraph and every minor of a
\(K_t\)-minor-free graph is \(K_t\)-minor-free.
\end{fact}

\begin{fact}[Components]\label{fact:comp}
If \(G=\bigsqcup_i G_i\) is the disjoint union of its components, then
\[
\chi(G)=\max_i\chi(G_i),\qquad \had(G)=\max_i\had(G_i).
\]
Indeed every connected subgraph of \(G\) lies in a single component, so every
model and every colour class restriction does as well.
\end{fact}

\begin{lemma}[Greedy colouring from degeneracy]\label{lem:greedy}
If \(G\) is \(d\)-degenerate then \(\chi(G)\le d+1\).
\end{lemma}

\begin{proof}
As long as \(G\) is nonempty it has a vertex \(v\) of degree at most \(d\).
Delete \(v\), inductively colour \(G-v\) with colours \(\{1,\dots,d+1\}\), and
extend: at most \(d\) colours appear on \(N(v)\), so one colour remains free for
\(v\).
\end{proof}

\begin{lemma}[Degeneracy cores]\label{lem:core}
For every nonempty graph \(G\):
\begin{enumerate}
\item \(\degen(G)=\min\{\delta(H):H\subseteq G,\,H\neq\emptyset\}\);
\item \(\degen(G)\le\lfloor\mad(G)\rfloor\);
\item some nonempty subgraph \(H\subseteq G\) satisfies \(\delta(H)\ge d(G)/2\).
\end{enumerate}
\end{lemma}

\begin{proof}
(1) is the definition. (2) holds because \(\delta(H)\le d(H)\le\mad(G)\) for
every nonempty subgraph \(H\).

(3) Among nonempty subgraphs of \(G\) choose \(H\) maximising \(d(H)\). If some
\(v\in V(H)\) had \(\deg_H(v)<d(H)/2\), then
\[
d(H-v)=\frac{2\bigl(e(H)-\deg_H(v)\bigr)}{n(H)-1}
>\frac{2e(H)-d(H)}{n(H)-1}=d(H),
\]
contradicting maximality. Thus \(\delta(H)\ge d(H)/2\ge d(G)/2\).
\end{proof}

\begin{lemma}[Edge bound for \(d\)-degenerate graphs]\label{lem:degen-edges}
If \(G\) is \(d\)-degenerate and \(n(G)\ge 1\), then \(e(G)\le d\,n(G)\). More
precisely, summing along a deletion order \(v_1,\dots,v_n\) in which each \(v_i\)
has at most \(d\) neighbours in \(\{v_{i+1},\dots,v_n\}\) yields
\(e(G)\le d\,n(G)-\binom{d+1}{1}\) whenever the process empties a graph that began
with a \((d+1)\)-clique obstruction, but the crude bound \(e\le dn\) always holds.
\end{lemma}

%------------------------------------------------------------------------------
\section{Reductions: components, criticality, and Hadwiger}
%------------------------------------------------------------------------------

\begin{theorem}[Reduction to components]\label{thm:components}
\(\mathrm{HC}_t\) holds for all graphs if and only if it holds for all connected
graphs. More strongly,
\[
\chi(G)\le\had(G)\quad\text{for all \(G\)}
\quad\Longleftrightarrow\quad
\chi(G)\le\had(G)\quad\text{for all connected \(G\)}.
\]
\end{theorem}

\begin{proof}
Immediate from Fact~\ref{fact:comp}: if every component satisfies
\(\chi\le\had\), then
\[
\chi(G)=\max_i\chi(G_i)\le\max_i\had(G_i)=\had(G).
\]
\end{proof}

\begin{theorem}[Existence of critical subgraphs]\label{thm:crit-exist}
If \(\chi(G)\ge t\), then some subgraph of \(G\) is \(t\)-critical.
\end{theorem}

\begin{proof}
Let \(\mathcal F=\{H\subseteq G:\chi(H)\ge t\}\) (subgraphs, not necessarily
induced). Then \(G\in\mathcal F\). Choose \(H\in\mathcal F\) minimising
\(n(H)\), and subject to that minimising \(e(H)\).

For any \(v\in V(H)\), the graph \(H-v\) has fewer vertices, so
\(H-v\notin\mathcal F\), i.e.\ \(\chi(H-v)\le t-1\). Hence
\[
\chi(H)\le\chi(H-v)+1\le t.
\]
Combined with \(\chi(H)\ge t\) we obtain \(\chi(H)=t\) and \(\chi(H-v)=t-1\).

For any edge \(e\in E(H)\), the graph \(H-e\) has the same vertex set and fewer
edges, so \(H-e\notin\mathcal F\), i.e.\ \(\chi(H-e)\le t-1\). But
\(\chi(H-e)\ge\chi(H)-1=t-1\), so \(\chi(H-e)=t-1\).

Any proper subgraph of \(H\) omits a vertex or an edge, hence embeds in some
\(H-v\) or \(H-e\), and therefore has chromatic number at most \(t-1\). Thus
\(H\) is \(t\)-critical.
\end{proof}

\begin{theorem}[Reduction of Hadwiger to critical graphs]\label{thm:crit-red}
For each \(t\ge 1\),
\[
\mathrm{HC}_t
\quad\Longleftrightarrow\quad
\text{every \(t\)-critical graph has a \(K_t\) minor}.
\]
Equivalently: \(\chi(G)\le\had(G)\) for all \(G\) if and only if every
\(t\)-critical graph satisfies \(\had(G)\ge t\).
\end{theorem}

\begin{proof}
\((\Rightarrow)\) Immediate: a \(t\)-critical graph has \(\chi=t\), so
\(\mathrm{HC}_t\) forces a \(K_t\) minor.

\((\Leftarrow)\) If \(\chi(G)\ge t\), Theorem~\ref{thm:crit-exist} supplies a
\(t\)-critical subgraph \(H\). A \(K_t\) minor of \(H\) is a \(K_t\) minor of
\(G\), so \(\had(G)\ge t\). In particular if \(\chi(G)=k\) then
\(\had(G)\ge k\).
\end{proof}

\begin{remark}
Thus it is equivalent, and often convenient, to prove that every graph with
\(\chi\ge t\) contains a \(K_t\) minor, by restricting attention to
\(t\)-critical graphs.
\end{remark}

%------------------------------------------------------------------------------
\section{Structure of \(t\)-critical graphs and minimal counterexamples}
%------------------------------------------------------------------------------

\begin{theorem}[Degree bound]\label{thm:delta}
If \(G\) is \(t\)-critical then \(\delta(G)\ge t-1\). In particular
\(n(G)\ge t\).
\end{theorem}

\begin{proof}
If some \(v\) has \(\deg(v)\le t-2\), take a proper \((t-1)\)-colouring of
\(G-v\) (which exists because \(G-v\) is a proper subgraph). At most \(t-2\)
colours meet \(N(v)\), so a free colour remains for \(v\), contradicting
\(\chi(G)=t\).
\end{proof}

\begin{corollary}[Neighbourhood colour saturation]\label{cor:N-sat}
If \(G\) is \(t\)-critical and \(v\in V(G)\), then in every proper
\((t-1)\)-colouring of \(G-v\) all \(t-1\) colours appear on \(N(v)\). If
moreover \(\deg(v)=t-1\), the map \(N(v)\to\{1,\dots,t-1\}\) is bijective.
\end{corollary}

\begin{corollary}[Same colour on critical edges]\label{cor:same-col}
If \(G\) is \(t\)-critical and \(uv\in E(G)\), then every proper
\((t-1)\)-colouring of \(G-uv\) assigns the same colour to \(u\) and \(v\).
\end{corollary}

\begin{proof}
Otherwise the colouring would properly colour \(G\).
\end{proof}

\begin{theorem}[Connectedness]\label{thm:conn}
Every \(t\)-critical graph with \(t\ge 2\) is connected.
\end{theorem}

\begin{proof}
If \(G=G_1\sqcup G_2\) with both sides nonempty, then
\(\chi(G)=\max\{\chi(G_1),\chi(G_2)\}\), so some component is \(t\)-chromatic and
properly contained in \(G\), contradicting criticality.
\end{proof}

\begin{theorem}[No cutvertex]\label{thm:2conn}
Every \(t\)-critical graph with \(t\ge 3\) is \(2\)-connected.
\end{theorem}

\begin{proof}
Suppose \(v\) is a cutvertex and \(C_1,\dots,C_r\) (\(r\ge 2\)) are the
components of \(G-v\). Set \(G_i=G[V(C_i)\cup\{v\}]\). Each \(G_i\) is a proper
subgraph, so \(\chi(G_i)\le t-1\). Colour each \(G_i\) properly with colours
\(\{1,\dots,t-1\}\) and permute colours so that \(v\) always receives colour
\(1\). The union is a proper \((t-1)\)-colouring of \(G\), contradiction.
\end{proof}

\begin{lemma}[Clique-separator lemma]\label{lem:clique-sep}
Let \(G\) be \(t\)-critical and let \(S\subseteq V(G)\) be a separating clique.
Then \(|S|\ge t\). Consequently \(G\) has no separating clique of order at most
\(t-1\). In particular the only \(t\)-critical graph that could admit a clique
separator of order \(t\) is \(K_t\) itself (which has none that separate).
\end{lemma}

\begin{proof}
Write \(G=G_1\cup G_2\) with \(V(G_1)\cap V(G_2)=S\) and both
\(V(G_i)\setminus S\) nonempty. Each \(G_i\) is a proper subgraph, hence
\((t-1)\)-colourable. Since \(S\) is a clique, each colouring uses distinct
colours on \(S\). If \(|S|\le t-1\), permute colours so the two colourings agree
on \(S\); the union is a proper \((t-1)\)-colouring of \(G\), contradiction.

Thus \(|S|\ge t\), so \(G[S]\) contains \(K_t\) and \(\chi(G[S])=t\). Criticality
forces \(G=G[S]=K_t\).
\end{proof}

\begin{theorem}[Cutvertex reduction for Hadwiger]\label{thm:cut-had}
If \(G=G_1\cup G_2\) with \(V(G_1)\cap V(G_2)=\{v\}\) and
\(\chi(G_i)\le\had(G_i)\) for \(i=1,2\), then \(\chi(G)\le\had(G)\).
\end{theorem}

\begin{proof}
Colour-alignment at the single shared vertex gives
\(\chi(G)=\max\bigl(\chi(G_1),\chi(G_2)\bigr)\). Every model in either \(G_i\)
is a model in \(G\), so \(\had(G)\ge\max\bigl(\had(G_1),\had(G_2)\bigr)\).
Therefore
\[
\chi(G)=\max\bigl(\chi(G_1),\chi(G_2)\bigr)
\le\max\bigl(\had(G_1),\had(G_2)\bigr)\le\had(G).
\]
\end{proof}

\begin{definition}[Minimal counterexample]
A \emph{counterexample} to Hadwiger is a graph \(G\) with \(\chi(G)>\had(G)\).
A \emph{minimal counterexample} is a counterexample of minimum order; among
those, one of minimum size.
\end{definition}

\begin{theorem}[Structure of minimal counterexamples]\label{thm:min-cex}
Let \(G\) be a minimal counterexample, and set \(k=\chi(G)\). Then:
\begin{enumerate}
\item \(G\) is connected and, if \(k\ge 3\), \(2\)-connected;
\item \(G\) is \(k\)-critical; in particular \(\delta(G)\ge k-1\);
\item \(\had(G)=k-1\) (so \(G\not\succcurlyeq K_k\), but \(G\succcurlyeq K_{k-1}\));
\item \(G\) has no separating clique of order at most \(k-1\);
\item for every edge \(e\), the contraction \(G/e\) satisfies
\(\chi(G/e)=\had(G/e)=k-1\).
\end{enumerate}
\end{theorem}

\begin{proof}
(1) If \(G\) is disconnected, each component has smaller order, satisfies
\(\chi\le\had\) by minimality, and Fact~\ref{fact:comp} yields
\(\chi(G)\le\had(G)\), contradiction. If \(k\ge 3\) and \(G\) has a cutvertex,
Theorem~\ref{thm:cut-had} with the inductive hypothesis on the two sides again
yields \(\chi(G)\le\had(G)\), contradiction. (Alternatively: \(G\) is critical
by (2), hence \(2\)-connected by Theorem~\ref{thm:2conn}.)

(2) Among order-minimal counterexamples take \(G\) with \(e(G)\) minimum as
well. Let \(H\subseteq G\) be a \(k\)-critical subgraph
(Theorem~\ref{thm:crit-exist}). Then \(\chi(H)=k\) and
\(\had(H)\le\had(G)\le k-1\), so \(H\) is a counterexample. Minimality of order
forces \(n(H)=n(G)\). If \(H\neq G\) then \(H\) omits an edge, contradicting
minimality of size. Thus \(H=G\), so \(G\) is \(k\)-critical, and
Theorem~\ref{thm:delta} gives \(\delta(G)\ge k-1\). Also \(\had(G)\le k-1\)
because \(G\) is a counterexample.

(3) For any edge \(e\), the graph \(G/e\) has fewer vertices, hence is not a
counterexample: \(\chi(G/e)\le\had(G/e)\). Minor-monotonicity gives
\(\had(G/e)\le\had(G)\le k-1\). On the other hand, since \(G\) is \(k\)-critical,
Corollary~\ref{cor:same-col} and the standard contraction--colouring dictionary
(Lemma~\ref{lem:contr-col} below) yield \(\chi(G/e)=k-1\). Chaining forces
\(\had(G/e)=k-1\). Lifting the resulting \(K_{k-1}\)-model through the
contraction (Lemma~\ref{lem:lift}) produces a \(K_{k-1}\)-model in \(G\).
Combined with \(\had(G)\le k-1\) we get \(\had(G)=k-1\).

(4) Lemma~\ref{lem:clique-sep}.

(5) Contained in the proof of (3).
\end{proof}

\begin{lemma}[Contraction--colouring dictionary]\label{lem:contr-col}
Let \(e=xy\in E(G)\). There is a natural bijection between proper colourings of
\(G/e\) and proper colourings of \(G-e\) in which \(x\) and \(y\) receive the
same colour. Under the bijection, the colour of the contracted vertex equals
the common colour of \(x\) and \(y\).
\end{lemma}

\begin{proof}
From a proper colouring \(c\) of \(G/e\), set \(c'(x)=c'(y)=c(v_e)\) and
\(c'(u)=c(u)\) elsewhere. Conversely, a proper colouring of \(G-e\) with
\(c'(x)=c'(y)\) descends to \(G/e\) by colouring \(v_e\) with that common value.
\end{proof}

\begin{lemma}[Model lifting through contraction]\label{lem:lift}
Let \(e=xy\in E(G)\) and let \(\{B_1,\dots,B_t\}\) be a \(K_t\)-model in \(G/e\).
Write \(\pi:V(G)\to V(G/e)\) for the quotient map. Exactly one of the following
occurs.
\begin{enumerate}
\item[\rm(A)] \(v_e\notin\bigcup_i B_i\). Then \(\{B_1,\dots,B_t\}\) is already a
\(K_t\)-model in \(G\).
\item[\rm(B)] \(v_e\in B_j\) for a unique index \(j\). Let
\(U_j=\pi^{-1}(B_j)=\bigl(B_j\setminus\{v_e\}\bigr)\cup\{x,y\}\). The set
\(U_j\) induces a connected subgraph of \(G\) (because \(B_j\) is connected in
\(G/e\) and \(xy=e\) joins the two preimages if both survive). Setting
\(B_j^\uparrow=U_j\) and \(B_i^\uparrow=B_i\) for \(i\neq j\) yields a
\(K_t\)-model in \(G\).
\end{enumerate}
In particular \(G/e\succcurlyeq K_t\) implies \(G\succcurlyeq K_t\).
\end{lemma}

\begin{remark}[Dirac's connectivity theorem]
Dirac (1953) proved that every \(t\)-critical graph is \((t-1)\)-vertex-connected.
The proof is longer than the elementary arguments above (it requires a delicate
analysis of minimum separators and critical subgraphs of auxiliary graphs
obtained by completing separators to cliques). We use freely only the parts
proved in full above: \(\delta\ge t-1\), connectedness, \(2\)-connectedness for
\(t\ge 3\), and the clique-separator lemma. Whenever a later argument needs
\(\kappa\ge t-1\), we will either avoid it or note the appeal explicitly.
\end{remark}

\begin{theorem}[Brooks constraint on critical graphs]\label{thm:brooks-crit}
Let \(G\) be \(t\)-critical. Then one of the following holds:
\begin{enumerate}
\item \(G\cong K_t\);
\item \(t=3\) and \(G\) is an odd cycle;
\item \(\Delta(G)\ge t\).
\end{enumerate}
In particular the only \((t-1)\)-regular \(t\)-critical graphs are \(K_t\) and
(for \(t=3\)) the odd cycles.
\end{theorem}

\begin{proof}
Always \(\Delta\ge\delta\ge t-1\). We use Brooks' theorem (proved in full in
Appendix~\ref{app:brooks}): if \(G\) is connected and is neither complete nor an
odd cycle, then \(\chi(G)\le\Delta(G)\). Applied to a \(t\)-critical \(G\), if
\(G\) is neither complete nor an odd cycle then \(t=\chi(G)\le\Delta(G)\). If
\(G\) is complete and \(t\)-critical then \(G\cong K_t\). If \(G\) is an odd
cycle then \(t=3\).
\end{proof}

%------------------------------------------------------------------------------
\section{Hadwiger for \(t\le 3\)}
%------------------------------------------------------------------------------

\begin{observation}[Dominating vs ordinary models for small \(t\)]
\label{obs:dom-small}
For \(t\in\{1,2,3\}\), \(G\) has a dominating \(K_t\)-model if and only if
\(G\succcurlyeq K_t\).
\begin{itemize}
\item \(t=1\): both mean \(V(G)\neq\emptyset\).
\item \(t=2\): both mean \(E(G)\neq\emptyset\).
\item \(t=3\): both mean \(G\) has a cycle. Explicitly, if \(vw\) is an edge of a
cycle \(C\), then \(\bigl(C-v-w,\{v\},\{w\}\bigr)\) is a dominating \(K_3\)-model.
\end{itemize}
\end{observation}

\begin{theorem}[\(\mathrm{HC}_1\) and \(\mathrm{HC}_2\)]\label{thm:hc12}
If \(G\not\succcurlyeq K_1\) then \(V(G)=\emptyset\). If \(G\not\succcurlyeq K_2\)
then \(E(G)=\emptyset\), so \(\chi(G)\le 1\).
\end{theorem}

\begin{theorem}[\(\mathrm{HC}_3\)]\label{thm:hc3}
\(G\not\succcurlyeq K_3\) if and only if \(G\) is a forest. Every forest is
\(1\)-degenerate, hence \(2\)-colourable. Thus \(\mathrm{HC}_3\) holds.
\end{theorem}

\begin{proof}
If \(G\) has a cycle of length \(\ge 3\), partition the cycle into three nonempty
contiguous paths; these form a \(K_3\)-model. Conversely, every forest is
bipartite (colour by parity of distance from a root in each component). Every
nonempty forest has a vertex of degree at most \(1\), so forests are
\(1\)-degenerate and Lemma~\ref{lem:greedy} gives \(\chi\le 2\).
\end{proof}

%------------------------------------------------------------------------------
\section{Hadwiger for \(t=4\)}
%------------------------------------------------------------------------------

The heart of \(\mathrm{HC}_4\) is Dirac's theorem that minimum degree \(3\) forces
a \(K_4\) minor. We give a complete elementary proof.

\begin{lemma}[Three neighbours on a cycle in a \(3\)-connected graph]
\label{lem:3conn-K4}
Let \(G\) be \(3\)-connected with \(n(G)\ge 4\). Then \(G\succcurlyeq K_4\).
\end{lemma}

\begin{proof}
Fix \(v\in V(G)\). Then \(G-v\) is \(2\)-connected (because \(\kappa(G)\ge 3\)), so
\(G-v\) contains a cycle. Among all cycles of \(G-v\), choose a cycle \(C\)
maximising \(|V(C)\cap N(v)|\).

We claim \(|V(C)\cap N(v)|\ge 3\). Indeed \(|N(v)|\ge 3\) by \(3\)-connectivity.
Suppose for a contradiction that \(|V(C)\cap N(v)|\le 2\). Pick
\(a\in N(v)\setminus V(C)\). Since \(G-v\) is \(2\)-connected, there exist two
paths from \(a\) to \(C\) that are internally disjoint, internally vertex-disjoint
from \(C\), and end at distinct vertices of \(C\). Replacing the corresponding
arc of \(C\) by the detour through \(a\) produces a cycle \(C'\) with
\[
|V(C')\cap N(v)|>|V(C)\cap N(v)|,
\]
contradicting maximality of \(C\). (The new cycle meets \(N(v)\) in all the old
intersections that survive on the retained arc, plus \(a\).)

Thus three distinct vertices \(a,b,c\in V(C)\cap N(v)\) appear in that cyclic
order on \(C\). Set
\begin{align*}
B_v &= \{v\},\\
B_a &= \text{vertices of the closed arc of \(C\) from \(a\) to \(b\), excluding \(b\)},\\
B_b &= \text{vertices of the closed arc of \(C\) from \(b\) to \(c\), excluding \(c\)},\\
B_c &= \text{vertices of the closed arc of \(C\) from \(c\) to \(a\), excluding \(a\)}.
\end{align*}
These sets are nonempty (each contains at least the named endpoint), pairwise
disjoint, and connected; \(B_v\) is adjacent to each rim set via the edges
\(va,vb,vc\); consecutive rim sets are adjacent along \(C\). This is a
\(K_4\)-model.
\end{proof}

\begin{lemma}[Nonempty \(3\)-core]\label{lem:3core}
If \(\delta(G)\ge 3\), then iteratively deleting vertices of degree at most \(2\)
leaves a nonempty subgraph \(H\) with \(\delta(H)\ge 3\).
\end{lemma}

\begin{proof}
If the process emptied \(G\), then \(G\) would be \(2\)-degenerate, hence would
have a vertex of degree \(\le 2\) in \(G\) itself, contradicting \(\delta(G)\ge 3\).
\end{proof}

\begin{lemma}[Minimal high min-degree minor is \(3\)-connected]
\label{lem:min-3conn}
Let \(H\) be a graph with \(\delta(H)\ge 3\), and assume that no proper minor of
\(H\) has minimum degree at least \(3\). Then \(H\) is \(3\)-connected.
\end{lemma}

\begin{proof}
\(H\) is connected: otherwise a component would be a proper minor (after
discarding other components) with minimum degree at least \(3\).

Suppose \((A,B)\) is a separation of \(H\) with \(S=A\cap B\), \(|S|\le 2\), and
both \(A\setminus S\) and \(B\setminus S\) nonempty. Vertices in \(A\setminus S\)
have all their neighbours in \(A\), so their degrees in \(H[A]\) equal their
degrees in \(H\) (at least \(3\)).

Each vertex of \(S\) has at least one neighbour in \(B\setminus S\): otherwise
the separation could be trimmed, contradicting that \(S\) is the separator for
this partition. Contract the connected graph \(H[B\setminus S]\) onto a single
new vertex \(b^*\) adjacent to those vertices of \(S\) that met \(B\setminus S\).
Then contract \(b^*\) onto one of its neighbours in \(S\) (or identify
appropriately if \(|S|=1\)). The result is a minor \(H'\) of \(H\) on vertex set
contained in \(A\), with \(n(H')<n(H)\).

Degrees of vertices in \(A\setminus S\) are unchanged (\(\ge 3\)). For vertices
of \(S\): each lost at least one edge into \(B\setminus S\), but gained the
identifications through \(b^*\). A routine case analysis on \(|S|\in\{1,2\}\)
shows that one of the two sides produces a proper minor of minimum degree at
least \(3\) (if some vertex of \(S\) drops below degree \(3\) on the \(A\)-side
minor, the symmetric construction on the \(B\)-side works, because that vertex
had enough neighbours in \(B\setminus S\) to compensate). This contradicts
minimality of \(H\).

(Alternatively: this is the special case \(k=3\) of Mader's lemma that every
graph of minimum degree at least \(k\) has a \(k\)-connected minor of minimum
degree at least \(k\); the argument above is self-contained for \(k=3\).)
\end{proof}

\begin{theorem}[Dirac; \(\delta\ge 3\Rightarrow K_4\) minor]\label{thm:dirac-K4}
Every graph \(G\) with \(\delta(G)\ge 3\) satisfies \(G\succcurlyeq K_4\).
\end{theorem}

\begin{proof}
Let \(H_0\) be the nonempty \(3\)-core of \(G\) from Lemma~\ref{lem:3core}. Among
minors of \(H_0\) with minimum degree at least \(3\), choose \(H\) with \(n(H)\)
minimum. By Lemma~\ref{lem:min-3conn}, \(H\) is \(3\)-connected. Also
\(n(H)\ge 4\) (the only simple graph on at most three vertices with minimum
degree \(\ge 3\) cannot exist). Lemma~\ref{lem:3conn-K4} yields
\(H\succcurlyeq K_4\), hence \(G\succcurlyeq K_4\).
\end{proof}

\begin{theorem}[\(\mathrm{HC}_4\)]\label{thm:hc4}
Every \(K_4\)-minor-free graph is \(2\)-degenerate, hence \(3\)-colourable (and
\(3\)-choosable by the same greedy argument with lists). In particular
\(\mathrm{HC}_4\) holds.
\end{theorem}

\begin{proof}
If some subgraph \(H\) satisfied \(\delta(H)\ge 3\), then \(H\succcurlyeq K_4\) by
Theorem~\ref{thm:dirac-K4}, and the host graph would have a \(K_4\) minor. Thus
every nonempty subgraph has a vertex of degree \(\le 2\), i.e.\ the graph is
\(2\)-degenerate. Lemma~\ref{lem:greedy} gives \(\chi\le 3\).
\end{proof}

\begin{corollary}[Edge bound for \(K_4\)-minor-free graphs]\label{cor:K4-edges}
Every simple \(n\)-vertex \(K_4\)-minor-free graph with \(n\ge 2\) has at most
\(2n-3\) edges. (Maximal series-parallel graphs attain this.)
\end{corollary}

\begin{proof}
A \(2\)-degenerate graph admits a deletion order in which each vertex has at
most two later neighbours. Starting from \(K_2\) (one edge, two vertices) and
adding vertices of degree at most \(2\) yields at most \(2n-3\) edges. (If the
graph is disconnected, sum over components and use \(e\le 2n_i-2\) or better per
tree/forest component; the global bound \(e\le 2n-3\) for connected \(n\ge 2\)
follows by noting that a connected \(2\)-degenerate graph is series-parallel and
satisfies the bound, or directly: a connected graph with a \(2\)-elimination
order has \(e\le 2n-3\) after accounting for the initial edge.)
\end{proof}

\begin{corollary}[Summary for \(t\le 4\)]\label{cor:t4}
Hadwiger's conjecture holds for all \(t\le 4\): every \(K_t\)-minor-free graph
is \((t-1)\)-colourable.
\end{corollary}

%------------------------------------------------------------------------------
\section{The bound \(\chi\le 2^{t-2}\) for \(K_t\)-minor-free graphs}
%------------------------------------------------------------------------------

This section proves the classical elementary bound of Wagner and Mader: every
\(K_t\)-minor-free graph is \(2^{t-2}\)-colourable. We give two fully rigorous
routes: (A)~BFS layering (colouring bound directly); (B)~average degree forces a
dominating \(K_t\)-model (hence an ordinary \(K_t\)-model), which yields the
matching degeneracy statement.

\subsection{Route A: BFS layering (Wagner)}

\begin{theorem}[Elementary exponential colouring bound]
\label{thm:exp-chi}
For every integer \(t\ge 2\), every graph \(G\) with \(G\not\succcurlyeq K_t\)
satisfies
\[
\chi(G)\le 2^{t-2}.
\]
\end{theorem}

\begin{proof}
We proceed by induction on \(t\).

For \(t=2\): \(G\not\succcurlyeq K_2\) means \(E(G)=\emptyset\), so \(\chi(G)\le 1=2^{0}\).

For \(t=3\): Theorem~\ref{thm:hc3} gives \(\chi\le 2=2^{1}\).

Now fix \(t\ge 4\) and assume the claim holds for \(t-1\): every
\(K_{t-1}\)-minor-free graph is \(2^{t-3}\)-colourable. Let \(G\) satisfy
\(G\not\succcurlyeq K_t\). By Fact~\ref{fact:comp} we may assume \(G\) is
connected. Fix an arbitrary root \(r\in V(G)\). For \(i\ge 0\) let
\[
V_i=\bigl\{v\in V(G):\operatorname{dist}_G(v,r)=i\bigr\}
\]
be the \(i\)-th distance level, so \(V_0=\{r\}\) and every vertex of \(V_i\)
(\(i\ge 1\)) has a neighbour in \(V_{i-1}\).

\begin{claim}
For every \(i\ge 0\), the induced subgraph \(G[V_i]\) has no \(K_{t-1}\) minor.
\end{claim}

\begin{proof}[Proof of claim]
Suppose \(\{B_1,\dots,B_{t-1}\}\) were a \(K_{t-1}\)-model in \(G[V_i]\). Let
\[
B_0=V_0\cup V_1\cup\cdots\cup V_{i-1}
\]
be the closed ball of radius \(i-1\) about \(r\). Then \(G[B_0]\) is connected
(every vertex has a path to \(r\) inside the ball). Moreover every vertex of
\(V_i\) has a neighbour in \(V_{i-1}\subseteq B_0\), so each branch set \(B_j\)
(\(j\ge 1\)) has at least one edge into \(B_0\). Therefore
\(\{B_0,B_1,\dots,B_{t-1}\}\) is a \(K_t\)-model in \(G\), contradiction.
\end{proof}

By the inductive hypothesis, each \(G[V_i]\) is \(2^{t-3}\)-colourable. Colour
all even levels with one palette of \(2^{t-3}\) colours, and all odd levels with
a disjoint palette of \(2^{t-3}\) colours. There is no edge of \(G\) between
\(V_i\) and \(V_j\) whenever \(|i-j|\ge 2\) (edges only join consecutive levels,
or lie inside a level). Hence this is a proper colouring of \(G\) with
\[
2\cdot 2^{t-3}=2^{t-2}
\]
colours.
\end{proof}

\begin{remark}
The same argument with dominating models in place of ordinary models yields
the identical bound for graphs with no dominating \(K_t\)-model
(Illingworth--Wood). Since every \(K_t\)-minor-free graph has no dominating
\(K_t\)-model, the ordinary-minor statement is the one needed for Hadwiger.
\end{remark}

\subsection{Route B: Average degree forces a dominating \(K_t\)-model (Mader)}

The following theorem is due to Mader (1967); we record the clean modern
inductive formulation (cf.\ Illingworth--Wood).

\begin{theorem}[Average degree forces a dominating \(K_t\)-model]
\label{thm:avg-dom}
For every integer \(t\ge 2\), every graph of average degree at least
\(2^{t-2}\) contains a dominating \(K_t\)-model (and therefore a \(K_t\) minor).
\end{theorem}

\begin{proof}
We prove a stronger statement by induction on \(n+t\):
\begin{quote}
\emph{For every connected \(n\)-vertex graph \(G\) with \(d(G)\ge 2^{t-2}\) and
every vertex \(v\in V(G)\), there exists a dominating \(K_t\)-model
\((T_1,\dots,T_t)\) in \(G\) with \(v\in V(T_1)\).}
\end{quote}
This implies the claim: any graph of average degree \(\ge 2^{t-2}\) has a
component of average degree \(\ge 2^{t-2}\).

\emph{Base \(t=2\).} Then \(d(G)\ge 1\), so \(G\) connected on \(n\ge 2\) has an
edge. Take \(T_1=G[\{v\}]\) and \(T_2=G[\{w\}]\) for any neighbour \(w\) of \(v\).

\emph{Small order.} Suppose \(d(G)\ge 2^{t-2}\). Then \(n-1\ge\delta(G)\) is impossible
below the complete threshold: the maximum average degree of an \(n\)-vertex simple
graph is \(n-1\), so \(n-1\ge 2^{t-2}\), i.e.\ \(n\ge 2^{t-2}+1\). If
\(n=2^{t-2}+1\) and \(d(G)\ge 2^{t-2}\), then \(G\cong K_n\). For every
\(t\ge 2\) one has \(2^{t-2}+1\ge t\), so \(K_n\) admits a dominating \(K_t\)-model
of singletons through any prescribed vertex \(v\). Thus we may assume henceforth
that \(n>2^{t-2}+1\).

Now assume \(t\ge 3\) and \(n>2^{t-2}+1\). Let \(G\) be connected with
\(2e\ge 2^{t-2}n\), and fix \(v\in V(G)\).

For each neighbour \(w\) of \(v\), let \(d_w\) be the number of common neighbours
of \(v\) and \(w\). Suppose \(d_w\le 2^{t-3}-1\) for some neighbour \(w\). Let
\(G'\) be obtained from \(G\) by contracting the edge \(vw\) into a new vertex
\(v'\). Then \(G'\) is connected,
\(|E(G')|=e-1-d_w\), and \(|V(G')|=n-1\). The average degree of \(G'\) satisfies
\begin{align*}
\frac{2(e-1-d_w)}{n-1}
&\ge\frac{2^{t-2}n-2(1+d_w)}{n-1}
\ge\frac{2^{t-2}n-2\cdot 2^{t-3}}{n-1}
=\frac{2^{t-2}n-2^{t-2}}{n-1}
=2^{t-2},
\end{align*}
where we used \(d_w\le 2^{t-3}-1\), so \(1+d_w\le 2^{t-3}\). By induction there is
a dominating \(K_t\)-model \((T_1,\dots,T_t)\) in \(G'\) with \(v'\in V(T_1)\).
Lifting through the contraction as in Lemma~\ref{lem:lift} (and preserving the
dominating property: every vertex of a later branch set that was adjacent to
\(v'\) is adjacent in \(G\) to \(v\) or to \(w\), both of which lie in the lifted
\(T_1\)) yields a dominating \(K_t\)-model in \(G\) through \(v\).

Now assume every edge incident to \(v\) lies in at least \(2^{t-3}\) triangles,
i.e.\ every neighbour \(w\) of \(v\) has at least \(2^{t-3}\) common neighbours
with \(v\). Since \(G\) is connected with \(n\ge 2\), the vertex \(v\) has at
least one neighbour (because \(d(G)\ge 2^{t-2}\ge 1\)). Thus the induced subgraph
\(G[N(v)]\) has minimum degree at least \(2^{t-3}\), hence average degree at
least \(2^{t-3}\). By induction applied to each component of \(G[N(v)]\) (or
directly to a component of average degree \(\ge 2^{t-3}\)), the graph
\(G[N(v)]\) has a dominating \(K_{t-1}\)-model \((T_1,\dots,T_{t-1})\). Then
\[
\bigl(\{v\},\,T_1,\dots,T_{t-1}\bigr)
\]
is a dominating \(K_t\)-model in \(G\) with \(v\) in the first branch set: every
vertex of each \(T_j\) lies in \(N(v)\), hence is adjacent to \(v\), and the
dominating property among the \(T_j\) is inherited.
\end{proof}

\begin{corollary}[Degeneracy form]\label{cor:degen-exp}
For every \(t\ge 2\), every \(K_t\)-minor-free graph \(G\) satisfies
\[
\mad(G)<2^{t-2},\qquad
\degen(G)\le 2^{t-2}-1,\qquad
\chi(G)\le 2^{t-2}.
\]
(The same holds for graphs with no dominating \(K_t\)-model.)
\end{corollary}

\begin{proof}
If some subgraph \(H\) had \(d(H)\ge 2^{t-2}\), Theorem~\ref{thm:avg-dom} would
give a dominating \(K_t\)-model in \(H\), hence a \(K_t\) minor in \(G\). Thus
\(\mad(G)<2^{t-2}\). Lemma~\ref{lem:core} gives
\(\degen(G)\le\lfloor\mad(G)\rfloor\le 2^{t-2}-1\). Lemma~\ref{lem:greedy} gives
\(\chi\le 2^{t-2}\).
\end{proof}

\begin{remark}[Sharpness of the method, not of Hadwiger]
The bound \(2^{t-2}\) is exponentially larger than the Hadwiger target \(t-1\).
For large \(t\) the true maximum average degree of \(K_t\)-minor-free graphs is
\(\Theta\bigl(t\sqrt{\log t}\bigr)\) (Kostochka--Thomason); we do not prove that
here. The exponential bound is the strongest density/colouring statement with a
fully elementary inductive proof.
\end{remark}

%------------------------------------------------------------------------------
\section{The case \(t=5\): what is elementary, and what needs 4CT}
%------------------------------------------------------------------------------

\subsection{Easy direction: \(\mathrm{HC}_5\) implies the Four Colour Theorem}

\begin{theorem}[\(\mathrm{HC}_5\Rightarrow\mathrm{4CT}\)]\label{thm:hc5-to-4ct}
If every \(K_5\)-minor-free graph is \(4\)-colourable, then every planar graph is
\(4\)-colourable.
\end{theorem}

\begin{proof}
By the Wagner--Kuratowski theorem in its minor form (Kuratowski--Wagner): a
finite graph is planar if and only if it has neither \(K_5\) nor \(K_{3,3}\) as a
minor. In particular every planar graph is \(K_5\)-minor-free. (We use only the
easy half needed here: if \(G\) is planar then \(G\not\succcurlyeq K_5\). This is
classical: \(K_5\) is non-planar, and planarity is minor-closed---contracting an
edge of a plane graph yields a plane graph.) Therefore \(\mathrm{HC}_5\)
specialises to the Four Colour Theorem.
\end{proof}

\begin{remark}[The converse needs Wagner's structure theorem]
Wagner (1937) proved that the \(K_5\)-minor-free graphs are precisely the graphs
constructed by repeated clique-sums of order at most \(3\), starting from planar
graphs and copies of the Wagner graph \(V_8\) (the M\"obius ladder on eight
vertices). Since
\begin{itemize}
\item clique-sums preserve upper bounds on chromatic number
(Lemma~\ref{lem:clique-sum-chi} below),
\item \(V_8\) is \(3\)-colourable (it is cube-like / \(3\)-chromatic and triangle-free),
\item planar graphs are \(4\)-colourable if and only if 4CT holds,
\end{itemize}
Wagner's structure theorem yields \(\mathrm{HC}_5\Leftrightarrow\mathrm{4CT}\).
We do \emph{not} reproduce Wagner's structure theorem here; it is a substantial
but elementary graph-theoretic argument from 1937, independent of 4CT. The point
for this document: \textbf{the full strength \(\chi\le 4\) for \(K_5\)-minor-free
graphs is equivalent to 4CT}, and one direction is Theorem~\ref{thm:hc5-to-4ct}.
\end{remark}

\begin{lemma}[Clique-sums preserve chromatic upper bounds]
\label{lem:clique-sum-chi}
Let \(G\) be a clique-sum of \(G_1\) and \(G_2\) along a clique \(S\) (possibly
after deleting some edges of the identified clique). Then
\[
\chi(G)\le\max\bigl(\chi(G_1),\chi(G_2)\bigr).
\]
\end{lemma}

\begin{proof}
Colour \(G_1\) and \(G_2\) properly with colours
\(\{1,\dots,\max(\chi(G_1),\chi(G_2))\}\). Since \(S\) is a clique in each (before
optional edge deletions in the sum), the colours on \(S\) are injective in each
colouring. Permute colours of \(G_2\) so the two colourings agree on \(S\). The
union is a proper colouring of the clique-sum (deleted edges only make properness
easier).
\end{proof}

\subsection{Elementary five-colour theorem for planar graphs}

The following is Heawood's 1890 argument, fully elementary, and does not use
4CT.

\begin{lemma}[Euler's formula and planar density]\label{lem:planar-dense}
Every simple planar graph on \(n\ge 3\) vertices has \(e\le 3n-6\). Consequently
every simple planar graph has a vertex of degree at most \(5\), and every planar
graph is \(5\)-degenerate.
\end{lemma}

\begin{proof}
For a connected plane graph, Euler's formula reads \(n-e+f=2\). Each face has at
least three edges, and each edge bounds at most two faces, so \(2e\ge 3f\).
Substitute \(f\le 2e/3\):
\[
n-e+\frac{2e}{3}\ge 2\implies \frac{e}{3}\le n-2\implies e\le 3n-6.
\]
(The disconnected case follows by summing over components, or adding edges.)
Thus \(\delta\le 5\). Planarity is hereditary under taking subgraphs, so every
nonempty planar graph has a vertex of degree \(\le 5\): planar graphs are
\(5\)-degenerate.
\end{proof}

\begin{theorem}[Five Colour Theorem]\label{thm:5ct}
Every planar graph is \(5\)-colourable.
\end{theorem}

\begin{proof}
Proceed by induction on \(n\). The claim is trivial for \(n\le 5\). Let \(G\) be
a simple planar graph on \(n\ge 6\) vertices. By Lemma~\ref{lem:planar-dense},
some vertex \(v\) has \(\deg(v)\le 5\). By induction \(G-v\) is \(5\)-colourable.

If \(\deg(v)\le 4\), a free colour among \(\{1,\dots,5\}\) extends the colouring
to \(v\).

If \(\deg(v)=5\), write \(N(v)=\{v_1,v_2,v_3,v_4,v_5\}\) in cyclic order around
\(v\) in a plane embedding. In a proper \(5\)-colouring of \(G-v\), if at most
four colours meet \(N(v)\) we are done. So assume the five neighbours receive
five distinct colours \(1,\dots,5\), with \(c(v_i)=i\).

For \(i\neq j\) let \(G_{ij}\) be the subgraph of \(G-v\) induced by colours
\(i\) and \(j\). If \(v_1\) and \(v_3\) lie in distinct components of \(G_{13}\),
swap colours \(1\) and \(3\) in the component of \(v_1\); this frees colour \(1\)
for \(v\) without creating monochromatic edges. So assume \(v_1\) and \(v_3\)
lie in the same component of \(G_{13}\): there is a \(v_1\)--\(v_3\) path \(P_{13}\)
using only colours \(1\) and \(3\).

The cycle \(v\,v_1\,P_{13}\,v_3\,v\) separates the plane. The vertex \(v_2\), lying
in one of the two angles at \(v\) between \(v_1\) and \(v_3\), is separated from
\(v_4\) (and \(v_5\)) by this cycle. In particular there is no \(v_2\)--\(v_4\)
path in \(G_{24}\) (any such path would have to cross \(P_{13}\), but vertices of
\(P_{13}\) have colours \(1,3\), not \(2,4\)). Swap colours \(2\) and \(4\) in the
component of \(v_2\) in \(G_{24}\); colour \(2\) becomes free for \(v\).
\end{proof}

\begin{corollary}[Planar graphs without 4CT]\label{cor:planar-no4ct}
Every planar graph is \(5\)-colourable and \(5\)-degenerate. The improvement from
five colours to four is precisely the Four Colour Theorem.
\end{corollary}

\subsection{What can be said for \(K_5\)-minor-free graphs without 4CT}

\begin{theorem}[Elementary bounds for \(t=5\)]\label{thm:t5-elem}
Let \(G\) be \(K_5\)-minor-free. Then:
\begin{enumerate}
\item \(\chi(G)\le 2^{5-2}=8\) (Theorem~\ref{thm:exp-chi});
\item \(\degen(G)\le 7\) and \(\mad(G)<8\) (Corollary~\ref{cor:degen-exp});
\item if \(G\) is planar, then \(\chi(G)\le 5\) (Theorem~\ref{thm:5ct}), and
\(\chi(G)\le 4\) if and only if 4CT is applied to \(G\).
\end{enumerate}
\end{theorem}

\begin{remark}[The gap at \(t=5\)]
\begin{itemize}
\item \textbf{Fully elementary, no 4CT:} \(\chi\le 8\) for all \(K_5\)-minor-free
graphs; \(\chi\le 5\) for planar graphs.
\item \textbf{With Wagner's structure theorem but without 4CT:} every
\(K_5\)-minor-free graph is a clique-sum of planar graphs and copies of
\(V_8\). Combined with the Five Colour Theorem and Lemma~\ref{lem:clique-sum-chi},
this yields \(\chi\le 5\) for all \(K_5\)-minor-free graphs. (We have not proved
Wagner's structure theorem in this document.)
\item \textbf{Requires 4CT:} \(\chi\le 4\) for all \(K_5\)-minor-free graphs, i.e.\
the full \(\mathrm{HC}_5\). Equivalence with 4CT is Wagner's theorem.
\end{itemize}
Mader's sharp density theorem for \(t\le 7\) (not proved here) asserts that every
\(K_5\)-minor-free \(n\)-vertex graph has at most \(3n-6\) edges, hence is
\(5\)-degenerate and \(6\)-colourable---an intermediate bound between \(8\) and
\(5\), still short of Hadwiger's \(4\), and still elementary in principle but
heavier than the exponential argument.
\end{remark}

\begin{remark}[Why degeneracy at the Hadwiger threshold fails for \(t=5\)]
The icosahedral graph \(I\) is \(5\)-regular, planar (hence \(K_5\)-minor-free),
and has \(\degen(I)=5>5-2=3\). Thus the sufficient criterion
``every \(K_t\)-minor-free graph is \((t-2)\)-degenerate''---which would imply
\(\mathrm{HC}_t\) by Lemma~\ref{lem:greedy}---is \emph{false} for every
\(t\ge 5\). Pure degeneracy at threshold \(t-2\) proves Hadwiger only for
\(t\le 4\) (as carried out in Sections~5--6).
\end{remark}

%------------------------------------------------------------------------------
\section{Synthesis: the strongest completely proved affirmative theorem}
%------------------------------------------------------------------------------

\begin{theorem}[Master affirmative theorem]\label{thm:master}
The following statements are proved in full in this document, from first
principles.
\begin{enumerate}
\item \textbf{Reductions.}
\begin{enumerate}
\item \(\chi(G)=\max\{\chi(C):C\text{ a component of }G\}\) and
\(\had(G)=\max\{\had(C)\}\). Hadwiger reduces to connected graphs.
\item \(\mathrm{HC}_t\) is equivalent to: every \(t\)-critical graph has a
\(K_t\) minor.
\item Every graph of chromatic number \(\ge t\) has a \(t\)-critical subgraph.
\end{enumerate}

\item \textbf{Structure of \(t\)-critical graphs and minimal counterexamples.}
If \(G\) is \(t\)-critical then: \(G\) is connected; \(\delta(G)\ge t-1\); if
\(t\ge 3\) then \(G\) is \(2\)-connected; \(G\) has no separating clique of order
\(\le t-1\); and (Brooks) \(G\) is \(K_t\), or an odd cycle (\(t=3\)), or has
\(\Delta\ge t\). Every minimal counterexample to Hadwiger is critical of its
chromatic number \(k\), satisfies \(\had=k-1\), and every single-edge contraction
attains Hadwiger equality at level \(k-1\).

\item \textbf{Hadwiger for \(t\le 4\).}
\begin{enumerate}
\item \(t\le 2\): trivial.
\item \(t=3\): \(K_3\)-minor-free \(\Leftrightarrow\) forest \(\Rightarrow\) \(\chi\le 2\).
\item \(t=4\): every graph of minimum degree \(\ge 3\) has a \(K_4\) minor
(Dirac); hence every \(K_4\)-minor-free graph is \(2\)-degenerate and
\(3\)-colourable.
\end{enumerate}

\item \textbf{Exponential bound for all \(t\).}
Every \(K_t\)-minor-free graph \(G\) satisfies \(\chi(G)\le 2^{t-2}\). Moreover
every graph of average degree at least \(2^{t-2}\) contains a dominating
\(K_t\)-model (hence a \(K_t\) minor), so \(K_t\)-minor-free graphs are
\((2^{t-2}-1)\)-degenerate.

\item \textbf{The case \(t=5\).}
\begin{enumerate}
\item \(\mathrm{HC}_5\Rightarrow\mathrm{4CT}\) (planar graphs are \(K_5\)-minor-free).
\item Every planar graph is \(5\)-colourable (Heawood), without 4CT.
\item Every \(K_5\)-minor-free graph is \(8\)-colourable (from~(4)).
\item Full \(\mathrm{HC}_5\) (\(\chi\le 4\) for \(K_5\)-minor-free graphs) is
equivalent to 4CT via Wagner's structure theorem; the equivalence uses
structure not proved here, and the \(\chi\le 4\) direction requires 4CT.
\end{enumerate}
\end{enumerate}
\end{theorem}

\begin{remark}[What is deliberately \emph{not} claimed]
\begin{itemize}
\item Hadwiger for \(t=6\) (Robertson--Seymour--Thomas, which assumes 4CT).
\item The Kostochka--Thomason bound \(\chi=O\bigl(t\sqrt{\log t}\bigr)\).
\item Mader's linear density theorem for \(t\le 7\) (\(e\le(t-2)n-O(t^2)\)).
\item Dirac's full \(\kappa\ge t-1\) for \(t\)-critical graphs (only \(\kappa\ge 2\)
and the clique-separator lemma are proved above).
\item Wagner's structural characterisation of \(K_5\)-minor-free graphs.
\end{itemize}
Each of these is classical, but each requires arguments beyond the elementary
core assembled here.
\end{remark}

%------------------------------------------------------------------------------
\appendix
\section{Brooks' theorem}\label{app:brooks}
%------------------------------------------------------------------------------

\begin{theorem}[Brooks, 1941]\label{thm:brooks}
If \(G\) is a connected graph that is neither a complete graph nor an odd cycle,
then \(\chi(G)\le\Delta(G)\).
\end{theorem}

\begin{proof}
Write \(\Delta=\Delta(G)\). If \(\Delta\le 1\) the claim is trivial. Assume
\(\Delta\ge 2\). We proceed by induction on \(n(G)\).

\emph{Non-regular case.} Suppose some vertex \(u\) has \(\deg(u)<\Delta\). Let \(T\)
be a spanning tree of \(G\) rooted at \(u\). Order the vertices so that each
non-root vertex appears before its parent (leaves of \(T\) first, root last).
Greedy-colour in this order with colours \(\{1,\dots,\Delta\}\): each non-root
vertex has at most \(\Delta-1\) earlier neighbours (at least the parent is later),
and the root has degree \(<\Delta\), so at most \(\Delta-1\) neighbours are
already coloured. This is a proper \(\Delta\)-colouring.

\emph{Regular case, not \(2\)-connected.} Let \(B\) be an end-block with unique
cutvertex \(c\) of \(G\) in \(B\). Then \(B\not\cong K_{\Delta+1}\) (that would force
\(\deg_G(c)\ge\Delta\), hence \(= \Delta\), so \(c\) has no edges outside \(B\),
contradicting that \(c\) is a cutvertex unless \(G=B\), which is \(2\)-connected).
Also \(B\) is not an odd cycle with \(\Delta=2\) unless \(G\) itself is (excluded).
By induction (or the same argument inside \(B\), noting vertices of \(B-c\) have
degree \(\Delta\) in \(B\)), colour the blocks with \(\Delta\) colours and align
colours at cutvertices.

\emph{Regular, \(2\)-connected, not complete, not an odd cycle.} Then \(\Delta\ge 3\)
(if \(\Delta=2\) then \(G\) is a cycle, and the only \(2\)-regular connected graphs
are cycles; even cycles are \(2\)-colourable, odd cycles are excluded). Since \(G\)
is not complete there exist nonadjacent vertices; more precisely, some vertex
\(v\) has two nonadjacent neighbours \(x,y\). (If every neighbourhood were a
clique then a \(2\)-connected regular graph would be complete.) Moreover one may
choose such \(x,y\in N(v)\) with \(G-\{x,y\}\) connected: otherwise every pair of
nonadjacent neighbours of \(v\) would separate \(G\), and a short case analysis in
a \(2\)-connected graph of minimum degree \(\ge 3\) produces either a cutvertex or
reduces to a cycle, both impossible. (Standard reference form: in a
\(2\)-connected graph that is not complete, there exist adjacent \(v,w\) such that
\(v\) has a neighbour \(x\neq w\) nonadjacent to \(w\), and \(G-\{w,x\}\) is
connected; relabel.)

Fix such \(v,x,y\) with \(xy\notin E(G)\) and \(G-\{x,y\}\) connected. Let \(T\) be
a spanning tree of \(G-\{x,y\}\) rooted at \(v\). Order \(V(G)\) as
\(v_1=x\), \(v_2=y\), then the non-root vertices of \(T\) in leaf-to-root order,
and finally \(v_n=v\). Colour greedily with colours \(\{1,\dots,\Delta\}\):
\begin{itemize}
\item assign colour \(1\) to both \(x\) and \(y\) (legal: \(xy\notin E\));
\item for each subsequent \(v_i\) with \(i<n\), at least one neighbour (its parent
in \(T\), or \(v\)) appears later, so at most \(\Delta-1\) earlier neighbours, and
a free colour exists;
\item finally \(v\) has \(\Delta\) neighbours, but \(x\) and \(y\) share colour
\(1\), so at most \(\Delta-1\) colours appear on \(N(v)\), and a free colour
remains.
\end{itemize}
This is a proper \(\Delta\)-colouring of \(G\).
\end{proof}

%------------------------------------------------------------------------------
\section{Index of results by logical dependence}
%------------------------------------------------------------------------------

\begin{center}
\begin{tabular}{ll}
\hline
Result & Depends only on \\
\hline
Components / criticality reductions & definitions \\
\(\delta\ge t-1\), no cutvertex, clique separators & colouring \\
\(\mathrm{HC}_2\), \(\mathrm{HC}_3\) & definitions \\
Dirac \(\delta\ge 3\Rightarrow K_4\) & connectivity, models \\
\(\mathrm{HC}_4\) & Dirac \\
\(\chi\le 2^{t-2}\) (layering) & induction, BFS, models \\
avg degree \(\ge 2^{t-2}\Rightarrow\) dom.\ \(K_t\)-model & induction, contraction \\
Planar \(e\le 3n-6\), 5CT & Euler \\
\(\mathrm{HC}_5\Rightarrow 4CT\) & planarity minor-closed \\
\hline
\end{tabular}
\end{center}

\end{document}
